\documentclass[11pt]{article}
\usepackage{amsmath}
\usepackage{graphicx}

\title{On the Slow Convergence of Impatient Methods}
\author{A.~Researcher \and B.~Coauthor}
\date{March 3, 2025}

\begin{document}
\maketitle

\begin{abstract}
  We revisit the classical fixed-point iteration and show that impatience costs
  precisely one order of convergence. The argument is elementary\footnote{It is
  elementary in the sense that it uses no machinery the reader does not already
  have.} and the constants are explicit.
\end{abstract}

\section{Introduction}
\label{sec:intro}

Fixed-point iteration is the oldest trick in numerical analysis~\cite{Ostrowski1966}.
Given a contraction $T$ on a complete metric space, the sequence $x_{n+1} = T(x_n)$
converges to the unique fixed point $x^\star$, and the rate is governed by the
Lipschitz constant $L < 1$:
\begin{equation}
  \label{eq:contraction}
  d(x_n, x^\star) \le \frac{L^n}{1 - L} \, d(x_1, x_0).
\end{equation}

What we add to this picture is the observation that \emph{stopping early} --- the
practice we call \textbf{impatience} --- changes the exponent rather than the
constant.

\section{Setting}

Throughout, $(X, d)$ is a complete metric space and $T \colon X \to X$ satisfies
$d(T(x), T(y)) \le L \, d(x,y)$ for all $x, y \in X$. We assume:
\begin{itemize}
  \item $0 < L < 1$, so that $T$ is a genuine contraction;
  \item $x_0 \in X$ is arbitrary, and $\alpha = d(x_1, x_0) > 0$;
  \item the tolerance $\varepsilon$ satisfies $\varepsilon < \alpha$.
\end{itemize}

The three assumptions are independent, and each is used exactly once below.

\subsection{The impatient stopping rule}

An \emph{impatient} run stops at the first index $n$ with
$d(x_{n+1}, x_n) \le \varepsilon$, rather than at the first $n$ with
$d(x_n, x^\star) \le \varepsilon$. The two differ by a factor of $1 - L$, which is
small exactly when the iteration is slow.

\begin{theorem}
  \label{thm:main}
  Let $n^\star$ be the impatient stopping index. Then
  $d(x_{n^\star}, x^\star) \le \varepsilon / (1 - L)$, and this bound is attained.
\end{theorem}

\begin{proof}
  Sum the geometric series $\sum_{k \ge 0} L^k$ from the triangle inequality, and
  note that equality holds for the affine map $T(x) = Lx$ on $\mathbb{R}$.
\end{proof}

\section{Experiments}

We ran the iteration for four values of $L$. The measured iteration counts are:

\begin{table}
  \begin{tabular}{lrr}
    \hline
    Contraction $L$ & Patient & Impatient \\
    \hline
    0.5  & 34  & 33 \\
    0.9  & 219 & 197 \\
    0.99 & 2291 & 1832 \\
    \hline
  \end{tabular}
  \caption{Iterations to reach $\varepsilon = 10^{-6}$.}
\end{table}

The gap widens as $L \to 1$, as Theorem~\ref{thm:main} predicts.

\begin{figure}
  \includegraphics[width=0.6\textwidth]{convergence.pdf}
  \caption{Distance to the fixed point against the iteration count, on a
  logarithmic scale.}
  \label{fig:convergence}
\end{figure}

\section{Conclusion}

Impatience is cheap when $L$ is small and expensive when it is not; a practitioner
who wants the bound of Equation~\eqref{eq:contraction} should divide the tolerance
by $1 - L$ before starting.\footnote{Or, equivalently, be patient.}

\end{document}
